\documentclass[runningheads]{llncs}

\usepackage{amsmath}
\usepackage{amssymb}
\usepackage{algorithm}
\usepackage{algpseudocode}
\usepackage{booktabs}
\usepackage{graphicx}
\usepackage{hyperref}
\usepackage{xspace}
\usepackage{microtype}

\renewcommand{\algorithmiccomment}[1]{\hfill\textit{\small// #1}}

\begin{document}

\title{Design of a Dual-Controller System for Safe Autonomous Remediation in Stateful and Stateless Workloads}

\titlerunning{FSM-Gated Dual-Track Autonomous Remediation}

\author{Pranav Negi\inst{1} \and Reema Sarkar\inst{1}}

\authorrunning{P. Negi, R. Sarkar}

\institute{Department of Computer Science and Engineering,
PES University\\
Supervised by Animesh Giri\\
\email{\{pes2ug23cs429,pes2ug23cs471\}@pesu.pes.edu}}

\maketitle

%------------------------------------------------------------------------
\begin{abstract}
Autonomous remediation of containerised workloads requires fundamentally
different safety guarantees depending on whether the workload carries
persistent state. Existing autonomic systems either treat all containers
uniformly or rely on coarse-grained health checks, leaving a gap between
stateless process-churn recovery and the careful, sequenced intervention
that stateful workloads demand. We present \textit{UnderCurrent}, a
dual-track autonomic control plane that runs two independent
remediation pipelines in parallel: a lightweight threshold-DAG controller
for stateless workloads (Type~S) and an FSM-gated, policy-enforced
controller for stateful workloads (Type~F). Both tracks attach live
eBPF kernel probes for sub-millisecond signal collection, compute
independent confidence scores over a sliding event window, and execute
every decision cycle on dual real and shadow paths to measure
behavioural divergence. A five-state formal FSM gates all Type~F
actions, preventing irreversible operations until a mandatory audit
phase completes. We evaluate UnderCurrent on a live Linux system with
real Docker containers across a nominal baseline run and a structured
fault-injection experiment. All six research metrics are computed
from real eBPF trace data. Results confirm zero misclassification,
perfect explanation completeness, full reversibility of dispatched
actions, and a divergence rate of 0.37 that directly quantifies the
safety contribution of accumulated FSM state history.
\end{abstract}

\keywords{autonomic computing \and eBPF \and container orchestration
\and finite state machines \and shadow execution \and safe remediation}

%------------------------------------------------------------------------
\section{Introduction}
\label{sec:intro}

Production containerised deployments routinely mix workloads with
fundamentally different failure semantics. A crashed \texttt{nginx}
worker is safely remediated by a restart; a postgres instance exhibiting
sustained high-latency block~I/O requires a carefully sequenced
flush-then-checkpoint procedure whose premature execution may corrupt
uncommitted transactions. Despite this distinction, contemporary
autonomic frameworks either apply a single restart-or-reschedule policy
uniformly across all containers~\cite{burnsk8s}, or delegate
remediation entirely to human operators who cannot respond at the
sub-second cadence that eBPF-level observability enables~\cite{bcc}.

The core difficulty is not signal collection---eBPF has made per-process
kernel telemetry cheap and always-on---but rather decision safety.
Autonomous remediation of stateful workloads requires three properties
that existing systems do not simultaneously provide: \textit{formal
sequencing} (an audit phase must precede any checkpoint operation),
\textit{reversibility enforcement} (irreversible actions must be
blocked until sequence constraints are satisfied), and
\textit{behavioural observability} (the system must expose, cycle by
cycle, whether accumulated history is materially changing decisions
relative to a stateless baseline).

We present \textit{UnderCurrent}, a dual-track control plane that
addresses these requirements. Our contributions are as follows.

\begin{enumerate}
\item A dual-track architecture that runs Type~S (stateless) and
Type~F (stateful) remediation pipelines independently on a shared
eBPF event bus, with per-track confidence scoring and sliding-window
state stores (\S\ref{sec:arch}).
\item A five-state formal FSM for Type~F that enforces action sequencing
and irreversibility constraints independently of the primary confidence
threshold DAG, evaluated via a policy enforcement layer that gates each
action at the execution boundary (\S\ref{sec:typef}).
\item A dual real/shadow execution architecture that runs every
reconcile cycle on both a live path and a stateless shadow path,
logging and quantifying divergence as a first-class metric (\S\ref{sec:arch}).
\item A six-metric evaluation framework with formal equations and
real-system results from live eBPF probes on Docker containers,
clearly distinguishing nominal-baseline from fault-injection findings
(\S\ref{sec:eval}).
\end{enumerate}

%------------------------------------------------------------------------
\section{Background and Related Work}
\label{sec:background}

\paragraph{Autonomic computing and MAPE-K.}
Kephart and Chess~\cite{mapek} established the MAPE-K loop---Monitor,
Analyse, Plan, Execute over a shared Knowledge base---as the canonical
architecture for self-managing systems. UnderCurrent maps directly onto
this model: eBPF probes constitute the Monitor layer, per-track
confidence scorers form the Analyse layer, the FSM-gated DAG is the
Plan layer, and the action executor plus policy sandbox implement the
Execute layer. The primary extension is the explicit split of Plan into
two independent branches (stateless DAG vs.\ stateful FSM-DAG) sharing
a common event bus.

\paragraph{Kubernetes operators and HPA/VPA.}
The Kubernetes Operator pattern~\cite{burnsk8s} and the Horizontal and
Vertical Pod Autoscalers provide declarative resource management, but
remediation actions are restricted to pod restarts, replica scaling, and
resource limit adjustments. Neither mechanism distinguishes stateful
from stateless workloads at the level of I/O-path signals, nor provides
formal sequencing constraints that prevent disruptive actions during
in-progress I/O. UnderCurrent complements operators by operating below
the Kubernetes API, directly on kernel-level observability.

\paragraph{eBPF observability.}
The BPF Compiler Collection (BCC)~\cite{bcc} and tools such as
bpftrace make it possible to attach probes to arbitrary kernel
functions with negligible runtime overhead. Projects such as
Cilium~\cite{cilium} apply eBPF to network policy enforcement.
UnderCurrent uses BCC to attach kprobes on \texttt{blk\_account\_io},
kretprobes on \texttt{vfs\_write} and \texttt{vfs\_read}, and
tracepoints on \texttt{syscalls/sys\_enter\_mount} and
\texttt{syscalls/sys\_exit\_mount}, giving direct visibility into I/O
latency, filesystem errors, and volume mount state without any
instrumentation of user-space application code.

\paragraph{Formal methods and runtime verification.}
TLA+~\cite{tlap} and Alloy provide specification languages for
verifying state-machine properties offline. Stateright~\cite{stateright}
enables runtime model checking. UnderCurrent takes a pragmatic middle
ground: the five-state FSM is implemented as executable Python with
transition guards, providing runtime enforcement of the same invariants
that a TLA+ specification would express. Formal verification of the FSM
transition table is identified as future work (\S\ref{sec:discussion}).

\paragraph{Shadow execution.}
Fowler's dark-launch pattern~\cite{fowler} and production shadowing
techniques run new code paths against live traffic without exposing
their output. UnderCurrent applies this concept to control-plane
decisions: every reconcile cycle executes on both a real path (decisions
executed, FSM mutated) and a shadow path (decisions computed against a
fresh throwaway FSM, never executed). Divergence between paths is logged
and quantified as metric M2.

\paragraph{Policy enforcement.}
Open Policy Agent (OPA)~\cite{opa} and Envoy's WASM extension
mechanism provide language-agnostic policy enforcement at service
boundaries. UnderCurrent's policy layer enforces FSM sequence
constraints at the action-execution boundary using a Python fallback
implementation, with optional upgrade to a WASM sandbox via
\texttt{wasmtime} for stronger isolation.

%------------------------------------------------------------------------
\section{System Architecture}
\label{sec:arch}

\subsection{Overview}

UnderCurrent runs two independent control-plane tracks in parallel
threads. Type~S handles stateless workloads; Type~F handles stateful
workloads. Both tracks share a common eBPF kernel probe infrastructure
and write decision traces to per-track \texttt{traces.jsonl} files.
Configuration parameters (sliding window length, confidence thresholds,
reconcile interval) are set at launch and shared across both tracks.

Figure~\ref{fig:arch} illustrates the architecture. Each track follows
the same pipeline: kernel probes emit events onto a per-track perf
buffer; events are consumed by a listener thread that filters by process
name and records them into a sliding-window StateStore; a reconcile
thread reads the StateStore on a fixed interval, computes a confidence
score, evaluates a decision DAG, and dispatches an action. Every
reconcile cycle is executed twice: once on the \textit{real path}
(live FSM, actions executed) and once on the \textit{shadow path}
(fresh throwaway FSM, actions suppressed). Divergences between the two
paths are logged to \texttt{divergence\_log.jsonl}.

\begin{figure}[t]
\centering
\fbox{\parbox{0.92\textwidth}{\small\centering
\textbf{UnderCurrent Dual-Track Architecture}\\[4pt]
\begin{tabular}{c}
Kernel eBPF Probes (process\_exit, tcp\_connect\_fail,\\
blk\_account\_io, vfs\_write, vfs\_read, sys\_enter/exit\_mount)\\[3pt]
\(\downarrow\) perf buffer per track\\[3pt]
\begin{tabular}{p{0.38\textwidth}|p{0.38\textwidth}}
\textbf{Type S (Stateless)} & \textbf{Type F (Stateful)}\\
EventListener + comm filter & EventListener + comm filter\\
StatelessStateStore (W=60s) & StatefulStateStore (W=60s)\\
Confidence Scorer S & Confidence Scorer F\\
Threshold DAG & FSM-Gated DAG\\
& Policy Enforcement Layer\\
\multicolumn{2}{c}{Action Executor}\\
\multicolumn{2}{c}{TraceLogger $\rightarrow$ traces.jsonl}\\
\end{tabular}\\[3pt]
\(\downarrow\) Real path + Shadow path (both tracks)\\
divergence\_log.jsonl
\end{tabular}
}}
\caption{UnderCurrent dual-track architecture. Both tracks attach eBPF
probes to the host kernel and filter events by container process name
before scoring and decision. Every cycle executes on both real and
shadow paths.}
\label{fig:arch}
\end{figure}

\subsection{eBPF Event Filtering}

All eBPF probes are kernel-global---they fire for every process on
the host. Events are attributed to Docker containers by reading the
kernel task's \texttt{comm} field via \texttt{bpf\_get\_current\_comm}
and matching it against a statically configured allow-list mapping
kernel process names to canonical container names (e.g.,
\texttt{redis-server}$\to$\texttt{redis}, \texttt{mysqld}$\to$\texttt{mysql}).
Events from processes not in the allow-list are discarded in the
userspace event handler before reaching the StateStore.

\subsection{Dual Real/Shadow Execution}

For each reconcile cycle, the track's \texttt{reconcile()} function
is called twice. The real invocation uses the live, persistent FSM
instance and dispatches decisions to the action executor. The shadow
invocation receives a freshly instantiated FSM with no accumulated
state; computed decisions are logged but never executed. Divergence is
detected by comparing the two decision sets and logging mismatches to
the divergence log. This architecture makes it possible to measure, in
every production cycle, how much accumulated FSM state history alters
outcomes relative to a memoryless baseline.

%------------------------------------------------------------------------
\section{Type~S: Stateless Control Plane}
\label{sec:types}

\subsection{Kernel Probes}

The Type~S track attaches two probes. A tracepoint on
\texttt{sched\_process\_exit} fires whenever any process exits,
capturing the process name and PID. A kprobe on \texttt{tcp\_connect}
fires at every outbound TCP connection attempt, capturing the initiating
process. Events are filtered to the configured container allow-list
before insertion into the StateStore.

\subsection{Confidence Scoring}

For each container $c$ and sliding window $W$, let $n_{\mathrm{pe}}(c,W)$
denote the count of \texttt{process\_exit} events and
$n_{\mathrm{tc}}(c,W)$ the count of \texttt{tcp\_connect\_fail} events
recorded in window $W$. The per-signal scores are:

\begin{equation}
s_{\mathrm{pe}}(c,W) = \min\!\left(\frac{n_{\mathrm{pe}}(c,W)}{\theta_{\mathrm{pe}}},\,1\right) \cdot \sigma_{\mathrm{pe}}
\label{eq:spe}
\end{equation}

\begin{equation}
s_{\mathrm{tc}}(c,W) = \min\!\left(\frac{n_{\mathrm{tc}}(c,W)}{\theta_{\mathrm{tc}}},\,1\right) \cdot \sigma_{\mathrm{tc}}
\label{eq:stc}
\end{equation}

where $\theta_{\mathrm{pe}}$, $\theta_{\mathrm{tc}}$ are scale
rates (events to saturation) and $\sigma_{\mathrm{pe}}$,
$\sigma_{\mathrm{tc}} < 1.0$ are per-signal score ceilings that bound
confidence below certainty. The combined confidence score is taken as
the maximum of the two signals:

\begin{equation}
\rho_S(c,W) = \max\!\left(s_{\mathrm{pe}}(c,W),\; s_{\mathrm{tc}}(c,W)\right)
\label{eq:rhos}
\end{equation}

The max combiner is chosen over a weighted sum to ensure that a single
dominant signal class drives the decision without being diluted by
orthogonal signals that happen to be quiet.

\subsection{Decision DAG}

The Type~S DAG maps confidence scores to three actions at two thresholds
$\tau_{\mathrm{restart}} = 0.80$ and $\tau_{\mathrm{reschedule}} = 0.95$:

\[
\delta_S(c) = \begin{cases}
\texttt{reschedule} & \text{if } \rho_S \geq \tau_{\mathrm{reschedule}}\\
\texttt{restart}    & \text{if } \tau_{\mathrm{restart}} \leq \rho_S < \tau_{\mathrm{reschedule}}\\
\texttt{no\_action} & \text{if } \rho_S < \tau_{\mathrm{restart}}
\end{cases}
\]

Both \texttt{restart} and \texttt{reschedule} are classified as
reversible actions. The Type~S track has no FSM and no policy
enforcement layer; the DAG is stateless by design, reflecting the
property that a fresh confidence snapshot is sufficient to determine
the appropriate response for a workload with no persistent on-disk state.

%------------------------------------------------------------------------
\section{Type~F: Stateful Control Plane}
\label{sec:typef}

\subsection{Kernel Probes}

The Type~F track attaches four probe groups. Kprobes on
\texttt{blk\_account\_io\_start} and \texttt{blk\_account\_io\_done}
bracket block-device I/O requests, recording the PID-keyed start
timestamp in a BPF hash map; at completion, the elapsed time is emitted
as a \texttt{blk\_io\_latency} event with a \texttt{latency\_us} field.
Kretprobes on \texttt{vfs\_write} and \texttt{vfs\_read} fire when
either syscall returns a negative value (indicating an error), emitting
\texttt{vfs\_write\_error} and \texttt{vfs\_read\_error} events
respectively. Tracepoints on \texttt{syscalls/sys\_enter\_umount} and
the exit of \texttt{sys\_enter\_mount} (filtered to successful returns)
emit \texttt{volume\_mount\_lost} and \texttt{volume\_mount\_restored}
events. A PID-keyed pending-mount hash map correlates mount entries with
their exit return codes.

\subsection{Confidence Scoring}

Let $n_h(c,W)$ be the count of \texttt{blk\_io\_latency} events for
container $c$ in window $W$ with $\texttt{latency\_us} \geq \Lambda$,
where $\Lambda = 10{,}000$\,\textmu s is the high-latency threshold.
Let $n_{\mathrm{vfs}}(c,W)$ be the total count of VFS error events.
Let $\mathbb{1}_{\mathrm{vol}}(c,W)$ be an indicator that evaluates
to 1 if any \texttt{volume\_mount\_lost} event in $W$ is not followed
by a subsequent \texttt{volume\_mount\_restored} (i.e., an unresolved
mount loss exists).

The three per-signal scores are:

\begin{equation}
s_{\mathrm{io}}(c,W) = \min\!\left(\frac{n_h(c,W)}{\theta_{\mathrm{io}}},\,1\right) \cdot \sigma_{\mathrm{io}}
\label{eq:sio}
\end{equation}

\begin{equation}
s_{\mathrm{vfs}}(c,W) = \min\!\left(\frac{n_{\mathrm{vfs}}(c,W)}{\theta_{\mathrm{vfs}}},\,1\right) \cdot \sigma_{\mathrm{vfs}}
\label{eq:svfs}
\end{equation}

\begin{equation}
s_{\mathrm{vol}}(c,W) = \sigma_{\mathrm{vol}} \cdot \mathbb{1}_{\mathrm{vol}}(c,W)
\label{eq:svol}
\end{equation}

with parameters $\theta_{\mathrm{io}} = 5$, $\sigma_{\mathrm{io}} = 0.85$,
$\theta_{\mathrm{vfs}} = 3$, $\sigma_{\mathrm{vfs}} = 0.88$, and
$\sigma_{\mathrm{vol}} = 0.97$. Volume mount loss receives the highest
possible score because an unresolved mount loss is an immediate data
availability threat requiring no accumulation period. The combined
Type~F confidence score is:

\begin{equation}
\rho_F(c,W) = \max\!\left(s_{\mathrm{io}}(c,W),\; s_{\mathrm{vfs}}(c,W),\; s_{\mathrm{vol}}(c,W)\right)
\label{eq:rhof}
\end{equation}

\subsection{Five-State FSM}

The Type~F control plane enforces a formal five-state FSM over the
lifecycle of each container's fault episode. Table~\ref{tab:fsm} gives
the complete transition relation. The FSM is persistent across reconcile
cycles; it is never reset unless a full Healthy recovery completes.

\begin{table}[t]
\caption{Type~F FSM transition relation. Each row gives a valid
(current state, trigger, next state) triple and the condition that
enables the trigger.}
\label{tab:fsm}
\centering
\begin{tabular}{llll}
\toprule
\textbf{From} & \textbf{Trigger} & \textbf{To} & \textbf{Condition}\\
\midrule
Healthy    & \texttt{degrade}         & Degraded   & $\rho_F \geq \tau_{\mathrm{flush}}$\\
Degraded   & \texttt{audit}           & Audited    & \texttt{flush\_io\_queue} executed\\
Audited    & \texttt{approve\_repair} & Repairing  & \texttt{checkpoint\_and\_restart} approved\\
Repairing  & \texttt{recover}         & Recovered  & action succeeded\\
Recovered  & \texttt{heal}            & Healthy    & $\rho_F < \tau_{\mathrm{flush}}$\\
Recovered  & \texttt{degrade}         & Degraded   & $\rho_F \geq \tau_{\mathrm{flush}}$ (re-fault)\\
\bottomrule
\end{tabular}
\end{table}

The mandatory Degraded$\to$Audited transition ensures that
\texttt{checkpoint\_and\_restart}---which interrupts active I/O and
may stall connections for several seconds---cannot fire until at least
one \texttt{flush\_io\_queue} cycle has confirmed that the container is
in a stable enough state to accept a restart. The \textit{production
lower bound on MTTR} is therefore at least one full reconcile interval
(default 5\,s) plus action execution time, regardless of confidence
score magnitude.

\subsection{FSM-Gated Reconciliation Pipeline}

Algorithm~\ref{alg:reconcile} details the FSM-gated reconcile
procedure. Steps are executed in order for each container in the
StateStore; the FSM instance is mutated in place on the real path
and used from a fresh copy on the shadow path.

\begin{algorithm}[t]
\caption{Type~F FSM-Gated Reconciliation (single container $c$)}
\label{alg:reconcile}
\begin{algorithmic}[1]
\Require StateStore $\mathcal{S}$, FSM instance $\mathcal{F}$, container $c$
\Ensure Decision record $d$ with action, FSM transition, blocked\_reason
\State $\rho \gets \rho_F(c, \mathcal{S})$ \Comment{Eq.~\ref{eq:rhof}}
\State $\mathit{sig} \gets \textsc{KernelSignals}(\mathcal{S}, c)$
\If{$\rho \geq \tau_{\mathrm{escalate}}$}
  \State $\mathit{raw} \gets \texttt{escalate}$
\ElsIf{$\rho \geq \tau_{\mathrm{repair}}$}
  \State $\mathit{raw} \gets \texttt{checkpoint\_and\_restart}$
\ElsIf{$\rho \geq \tau_{\mathrm{flush}}$}
  \State $\mathit{raw} \gets \texttt{flush\_io\_queue}$
\Else
  \State $\mathit{raw} \gets \texttt{no\_action}$
\EndIf
\State $\mathit{action} \gets \mathit{raw}$
\If{$\mathit{raw} \in \mathit{IrreversibleActions}$} \Comment{Reversibility gate}
  \State $\mathit{action} \gets \texttt{no\_action}$; set \textit{blocked\_reason}
\EndIf
\If{$\rho \geq \tau_{\mathrm{flush}}$ \textbf{and} $\mathit{action} \neq \texttt{no\_action}$} \Comment{FSM degrade}
  \If{$\mathcal{F}.\mathit{state}(c) \in \{\textit{Healthy, Recovered}\}$}
    \State $\mathcal{F}.\mathit{apply}(c, \texttt{degrade})$
  \EndIf
\EndIf
\If{$\mathit{action} = \texttt{checkpoint\_and\_restart}$ \textbf{and} $\mathcal{F}.\mathit{state}(c) \neq \textit{Audited}$}
  \State $\mathit{action} \gets \texttt{flush\_io\_queue}$; set \textit{blocked\_reason} \Comment{FSM gate}
\EndIf
\If{$\mathit{action} = \texttt{flush\_io\_queue}$ \textbf{and} $\mathcal{F}.\mathit{state}(c) = \textit{Degraded}$}
  \State $\mathcal{F}.\mathit{apply}(c, \texttt{audit})$ \Comment{Degraded $\to$ Audited}
\EndIf
\If{$\mathit{action} = \texttt{checkpoint\_and\_restart}$ \textbf{and} $\mathcal{F}.\mathit{state}(c) = \textit{Audited}$}
  \State $\mathcal{F}.\mathit{apply}(c, \texttt{approve\_repair})$ \Comment{Audited $\to$ Repairing}
\EndIf
\If{$\rho < \tau_{\mathrm{flush}}$ \textbf{and} $\mathcal{F}.\mathit{state}(c) = \textit{Recovered}$}
  \State $\mathcal{F}.\mathit{apply}(c, \texttt{heal})$ \Comment{Recovered $\to$ Healthy}
\EndIf
\State \Return decision record $(c, \rho, \mathit{action}, \mathit{sig},
  \mathcal{F}.\mathit{state}(c), \mathit{blocked\_reason})$
\end{algorithmic}
\end{algorithm}

Threshold values are $\tau_{\mathrm{flush}} = 0.50$,
$\tau_{\mathrm{repair}} = 0.80$, and $\tau_{\mathrm{escalate}} = 0.95$.

\subsection{Policy Enforcement Layer}

A pluggable policy enforcement layer sits between Algorithm~\ref{alg:reconcile}
and the action executor. It re-evaluates FSM sequence constraints
independently of the DAG, providing a secondary enforcement boundary
that is decoupled from the primary reconcile logic. The layer is
implemented as a Python module with a defined interface that accepts
a decision record and returns an allow/deny verdict. An optional WASM
sandbox via \texttt{wasmtime} provides stronger isolation by executing
the policy in a restricted execution environment; if \texttt{wasmtime}
is not installed the system falls back to the Python implementation
and logs a diagnostic warning. In our evaluation, the Python fallback
was active; WASM isolation was not exercised.

%------------------------------------------------------------------------
\section{Research Metrics and Evaluation Framework}
\label{sec:metrics}

We define six metrics over the set $\mathcal{T}$ of all trace records
produced by the system. Let $\mathcal{T}_{\mathrm{real}} \subseteq
\mathcal{T}$ denote records written on the real execution path
($\texttt{mode} = \texttt{real}$), and let
$\mathcal{T}_{\mathrm{active}} = \{t \in \mathcal{T}_{\mathrm{real}} :
\texttt{action}(t) \neq \texttt{no\_action}\}$ denote real-mode records
where an active decision was taken.

\paragraph{M1: Misclassification Rate.}
A trace record $t$ is \textit{misclassified} if it carries signals
inconsistent with its \texttt{node\_type} label---for example, a
\texttt{blk\_io\_latency} signal in a Type~S record, or a
\texttt{process\_exit} signal in a Type~F record.

\begin{equation}
M_1 = \frac{|\{t \in \mathcal{T} : \mathrm{misclassified}(t)\}|}{|\mathcal{T}|}
\label{eq:m1}
\end{equation}

\paragraph{M2: Divergence Rate.}
Let $a_{\mathrm{real}}(t)$ and $a_{\mathrm{shadow}}(t)$ denote the
actions chosen on the real and shadow paths for the same container and
cycle. A divergence occurs when the two paths disagree.

\begin{equation}
M_2 = \frac{|\{t \in \mathcal{T}_{\mathrm{real}} : a_{\mathrm{real}}(t) \neq a_{\mathrm{shadow}}(t)\}|}{|\mathcal{T}_{\mathrm{real}}|}
\label{eq:m2}
\end{equation}

\paragraph{M3: Mean Time To Recovery (MTTR).}
Let $\mathcal{E}$ be the set of complete fault episodes, where each
episode $e \in \mathcal{E}$ is a (container, time-ordered) pair of
a first non-\texttt{no\_action} decision at $t_{\mathrm{onset}}(e)$
and a first subsequent \texttt{no\_action} decision at
$t_{\mathrm{recovery}}(e)$, both timestamped by
\texttt{decision\_timestamp}.

\begin{equation}
M_3 = \frac{1}{|\mathcal{E}|}\sum_{e \in \mathcal{E}} \bigl(t_{\mathrm{recovery}}(e) - t_{\mathrm{onset}}(e)\bigr)
\label{eq:m3}
\end{equation}

\paragraph{M4: Reversibility Ratio.}
Let $\mathrm{reversible}(t) = 1$ if the action in record $t$ is
classified as reversible in the action registry.

\begin{equation}
M_4 = \frac{|\{t \in \mathcal{T}_{\mathrm{active}} : \mathrm{reversible}(t)\}|}{|\mathcal{T}_{\mathrm{active}}|}
\label{eq:m4}
\end{equation}

\paragraph{M5: Explanation Completeness.}
A trace record $t$ is \textit{complete} if the fields
\texttt{why}, \texttt{action}, \texttt{score}, \texttt{kernel\_signals},
and \texttt{dag\_pattern} are all non-empty.

\begin{equation}
M_5 = \frac{|\{t \in \mathcal{T} : \mathrm{complete}(t)\}|}{|\mathcal{T}|}
\label{eq:m5}
\end{equation}

\paragraph{M6: Unsafe-Action Suppression Rate.}
A real-mode record is \textit{suppressed} if the FSM gate or policy
enforcement layer changed the action from the raw DAG output
(i.e., $\texttt{blocked\_reason}(t) \neq \varnothing$).

\begin{equation}
M_6 = \frac{|\{t \in \mathcal{T}_{\mathrm{real}} : \mathrm{suppressed}(t)\}|}{|\mathcal{T}_{\mathrm{real}}|}
\label{eq:m6}
\end{equation}

\subsection{Cross-Controller Scenario Table}

Table~\ref{tab:scenarios} summarises the functional correctness
scenarios used to validate decision accuracy across both tracks.
Scenarios exercise sub-threshold, minor, major, and critical severity
levels, plus a fault-then-recovery transition.

\begin{table}[t]
\caption{Cross-controller functional correctness scenarios. Score
column shows input confidence score; Expected and Observed columns
give the expected and observed action. All scenarios passed.}
\label{tab:scenarios}
\centering
\begin{tabular}{lllll}
\toprule
\textbf{Track} & \textbf{Scenario} & \textbf{Score} & \textbf{Expected} & \textbf{Observed}\\
\midrule
S & Sub-threshold     & 0.54 & \texttt{no\_action}  & \texttt{no\_action}\\
S & Minor fault       & 0.90 & \texttt{restart}     & \texttt{restart}\\
S & Critical fault    & 0.95 & \texttt{reschedule}  & \texttt{reschedule}\\
S & Recovery          & 0.95$\to$drop & \texttt{reschedule} & \texttt{reschedule}\\
F & Sub-threshold     & 0.29 & \texttt{no\_action}  & \texttt{no\_action}\\
F & Minor (flush)     & 0.59 & \texttt{flush\_io\_queue} & \texttt{flush\_io\_queue}\\
F & Major (FSM-gated) & 0.85 & \texttt{flush\_io\_queue}\textsuperscript{*} & \texttt{flush\_io\_queue}\\
F & Critical          & 0.97 & \texttt{escalate}    & \texttt{escalate}\\
F & Recovery          & 0.59$\to$drop & \texttt{flush\_io\_queue} & \texttt{flush\_io\_queue}\\
\bottomrule
\multicolumn{5}{l}{\textsuperscript{*}FSM in Degraded (not Audited); \texttt{checkpoint\_and\_restart}
deferred to \texttt{flush\_io\_queue}.}
\end{tabular}
\end{table}

%------------------------------------------------------------------------
\section{Evaluation}
\label{sec:eval}

All results reported in this section are computed from real-system eBPF
trace data collected on a live Linux host. No simulation or synthetic
data generation was used.

\subsection{Experimental Setup}

The host is an x86-64 developer workstation running Ubuntu 22.04 with
Linux kernel 5.15, a 64\,GB virtual disk backed by a rotating hard
drive (HDD), and 8\,GB RAM. UnderCurrent was deployed directly on the
host (not inside Kubernetes) using BCC for eBPF probe attachment,
requiring \texttt{CAP\_BPF} / root privileges. Three Docker containers
were active throughout both evaluation phases: \texttt{nginx} (reverse
proxy), \texttt{postgres} (relational database with a persistent volume
mount), and \texttt{redis} (in-memory store). MySQL was not available
at evaluation time. The reconcile interval was 5\,s on both tracks;
the sliding event window was 60\,s.

Two evaluation phases were conducted. Phase~1 was a 32-minute nominal
baseline run with no induced faults, providing evidence of correct
behaviour under healthy conditions. Phase~2 was a structured
fault-injection run of approximately 5.5 minutes in which VFS write
errors were induced inside the \texttt{postgres} container by
repeatedly executing \texttt{COPY ... TO '/dev/full'} via
\texttt{psql}, and block~I/O stress was applied by a parallel
\texttt{dd conv=fdatasync} loop overwriting a fixed 1\,MB temporary
file.

Metric values were computed by \texttt{shared/real\_metrics\_report.py},
which filters all inputs to \texttt{mode=real} entries before any
calculation and uses \texttt{decision\_timestamp} (set at decision
time, before action execution) rather than \texttt{trace\_time} (set
after \texttt{execute()} returns) for M3, avoiding inflation from
action execution latency.

\subsection{Phase~1: Nominal Baseline}

Over 32 minutes, the Type~S track processed approximately 395 reconcile
cycles across nginx, postgres, and redis, recording 1{,}187 real-mode
stateless trace entries. The Type~F track processed 386 cycles with 386
real-mode stateful trace entries, all for postgres (the only container
generating qualifying block-device I/O signals). The stateful FSM
remained in the Healthy state for every one of the 386 cycles, issuing
\texttt{no\_action} throughout.

The absence of stateful fault episodes in the nominal run is not a null
result. It confirms that UnderCurrent does not over-trigger on healthy
containers, which is a critical safety property for a system that
governs irreversible remediation actions. A system that issued spurious
\texttt{checkpoint\_and\_restart} decisions during normal operation
would introduce availability risk with no fault-mitigation benefit.
The 307 captured \texttt{blk\_io\_latency\_normal} events demonstrate
that block-I/O telemetry was active and reaching the StateStore;
all observed latencies fell below the $\Lambda = 10{,}000$\,\textmu s
threshold.

\subsection{Phase~2: Fault Injection}

At the time of metric collection (following the 5.5-minute fault
injection and a 2-minute cooldown), the Type~F track had recorded 119
real-mode trace entries and 44 divergence log entries; the Type~S track
had recorded 123 real-mode entries.

Table~\ref{tab:results} summarises all six metric values from both
evaluation phases.

\begin{table}[t]
\caption{Metric values from real-system eBPF evaluation. Phase~1
results use data from the 32-minute nominal run ($N_S=1{,}187$,
$N_F=386$ real entries). Phase~2 results use data from the fault
injection run ($N_F=119$ real entries). All values computed by
\texttt{real\_metrics\_report.py} from \texttt{mode=real} trace
entries only.}
\label{tab:results}
\centering
\begin{tabular}{lllll}
\toprule
\textbf{Metric} & \textbf{Type S} & \textbf{Type F (Ph.1)} & \textbf{Type F (Ph.2)} & \textbf{Note}\\
\midrule
M1 (misclassification)    & 0.0000 & 0.0000 & 0.0000 & $N=3{,}146$ combined\\
M2 (divergence rate)      & 0.0000 & 0.0000 & 0.3697 & 44 of 119 cycles\\
M3 (MTTR, seconds)        & ---    & ---    & 242.8  & $N=2$ episodes\textsuperscript{a}\\
M4 (reversibility)        & 1.0000 & ---    & 1.0000 & $N=1{,}039$ / $N=95$\\
M5 (explanation complete) & 1.0000 & 1.0000 & 1.0000 & Zero missing fields\\
M6 (suppression rate)     & 0.0000 & 0.0000 & 0.3697 & 44 FSM-gated actions\textsuperscript{b}\\
\bottomrule
\multicolumn{5}{l}{\textsuperscript{a}Episodes: 61.0\,s and 424.6\,s; mean = 242.8\,s; $N=2$ is statistically thin.}\\
\multicolumn{5}{l}{\textsuperscript{b}Policy layer ran in Python fallback mode; WASM sandbox not exercised.}
\end{tabular}
\end{table}

\textbf{M1 = 0.0000.} No trace record carried signals inconsistent
with its track label across the combined 3{,}146 traces from both
phases. This confirms that the per-track eBPF probes and the
\texttt{node\_type} label assignment are consistent throughout.

\textbf{M4 = 1.0000 on both tracks.} Every one of the 1{,}039
Type~S active decisions dispatched during Phase~1 was a
\texttt{reschedule} or \texttt{restart} action, both classified as
reversible. During Phase~2, all 95 Type~F active decisions were
\texttt{flush\_io\_queue} or \texttt{checkpoint\_and\_restart};
both are classified as reversible in the action registry. The FSM gate
successfully prevented any irreversible action from reaching the
executor in either phase.

\textbf{M5 = 1.0000.} All trace records on both tracks in both phases
carried complete explanatory fields (\texttt{why}, \texttt{action},
\texttt{score}, \texttt{kernel\_signals}, \texttt{dag\_pattern}). Zero
records had missing fields. This confirms that the trace schema is
populated consistently regardless of execution path (real or shadow)
or fault state.

\textbf{M2 = M6 = 0.0000 (Phase~1).} With all containers healthy and
no fault signals crossing any threshold, the real and shadow paths
produced identical decisions on every cycle. The FSM gate had no effect
because no raw DAG action exceeded \texttt{no\_action}.

\textbf{M2 = 0.3697 (Phase~2).} During fault injection, 44 of 119
real-mode stateful cycles diverged between the real and shadow paths.
All 44 divergences had the pair $(\texttt{checkpoint\_and\_restart},
\texttt{flush\_io\_queue})$, meaning the real path chose
\texttt{checkpoint\_and\_restart} while the shadow path chose
\texttt{flush\_io\_queue}. This pattern arises because the real FSM,
having accumulated state across multiple cycles, had reached the Audited
state and approved the repair action; the shadow FSM, freshly
instantiated at the start of each cycle, could advance at most to
Degraded within a single cycle, causing the FSM gate to defer
\texttt{checkpoint\_and\_restart} to \texttt{flush\_io\_queue}.

\textbf{M6 = 0.3697 (Phase~2).} Exactly 44 real-mode records carried a
non-null \texttt{blocked\_reason}, indicating that the FSM gate or
policy layer changed the action from the raw DAG output. These 44
suppressions correspond to cycles in which the raw DAG called for
\texttt{checkpoint\_and\_restart} (score $\geq \tau_{\mathrm{repair}}$)
but the real FSM was in a state that did not yet satisfy the Audited
precondition (typically Degraded, immediately following re-entry from
a Recovered state after a prior successful checkpoint cycle). The
policy enforcement layer (Python fallback) confirmed the FSM verdict
in all cases.

\textbf{M3 = 242.8\,s (Phase~2, N=2).} Two complete fault episodes
were observed for postgres: one episode of 61.0\,s and one of 424.6\,s,
yielding a mean of 242.8\,s. The wide spread reflects the different
fault intensities of the two episodes rather than system non-determinism.
The $N=2$ sample is statistically thin; this value is directionally
valid but should be treated with caution. The production lower bound on
MTTR for Type~F is at least one full reconcile interval (5\,s) because
the mandatory Degraded$\to$Audited transition requires at least one
\texttt{flush\_io\_queue} cycle before \texttt{checkpoint\_and\_restart}
can be approved.

\subsection{Cross-Controller Functional Correctness}

All nine scenarios in Table~\ref{tab:scenarios} produced the expected
action, giving a cross-controller decision accuracy of 1.00 at all
tested severity levels. Notably, the FSM-gated major-fault scenario
(score $= 0.85$, FSM in Degraded) correctly deferred
\texttt{checkpoint\_and\_restart} to \texttt{flush\_io\_queue},
confirming that the FSM gate operates independently of the raw DAG
output.

\subsection{The M2 = M6 Equality and Its Architectural Meaning}

The numerical equality of M2 and M6 in Phase~2 is architecturally
significant and warrants explicit discussion. Both metrics equal
$44/119 = 0.3697$ because the 44 cycles on which the shadow path
diverged from the real path are precisely the cycles on which the real
FSM had accumulated enough state to approve \texttt{checkpoint\_and\_restart},
while in a symmetric set of 44 cycles the real FSM gate suppressed
\texttt{checkpoint\_and\_restart} (an M6 event) during the
re-Degraded phase before reaching Audited.

More precisely: the fault injection induced a recurring
fault-recovery-re-fault cycle in which the FSM alternated between
Audited (allowing C\&R, causing shadow divergence because the shadow
cannot reach Audited in one cycle) and Degraded (suppressing C\&R,
causing M6 events because the raw DAG still scored $\geq 0.80$). The
equal count arises from the regularity of this pattern under the
specific fault workload.

The structural relationship between M2 and M6 is that both are driven
by the same underlying mechanism: the difference between accumulated
FSM state (real path) and memoryless FSM state (shadow path / raw DAG).
M2 quantifies how often that difference changes outcomes; M6 quantifies
how often the safety mechanism acts on that difference. In any
deployment where M2 $>$ 0, FSM history is materially influencing
decisions; in any deployment where M6 $>$ 0, the safety gate is
actively doing work. Their equality in this evaluation reflects the
specific fault-recovery cadence and is not an identity that holds in
general.

%------------------------------------------------------------------------
\section{Discussion}
\label{sec:discussion}

\subsection{Design Choices}

The max combiner in Eqs.~\ref{eq:rhos} and~\ref{eq:rhof} was chosen
over a weighted sum to prevent a quiescent signal class from diluting
the dominant fault indicator. A container with frequent process exits
but no TCP failures would receive the correct high score from $s_{pe}$
regardless of whether $s_{tc}$ is zero. The sub-unity score ceilings
($\sigma < 1.0$) reserve the top of the $[0,1]$ range for future
signal classes with higher diagnostic confidence, such as hardware
error counters or out-of-memory kernel events.

The choice to make the Type~S DAG stateless (no FSM) is deliberate.
Stateless workloads, by definition, carry no on-disk state that could
be corrupted by an untimely restart; the risk profile of
\texttt{restart} and \texttt{reschedule} is symmetric between a fresh
container and a running one. Adding FSM sequencing to the stateless
track would impose audit-phase latency without a corresponding safety
benefit.

The dual real/shadow architecture imposes a per-cycle overhead of one
additional DAG evaluation and FSM instantiation. In our implementation
the shadow computation completes within a few microseconds because it
operates on in-memory data structures. The architectural benefit---a
live, cycle-by-cycle measurement of how much FSM history is changing
outcomes---justifies this overhead for any deployment in which operator
trust in the autonomous system is still being established.

\subsection{Limitations}

Several limitations should be noted. First, this evaluation was
conducted on a developer workstation, not a production Kubernetes
cluster. The eBPF probe architecture is kernel-version-independent
above Linux 4.9, but the specific kprobe attachment points
(\texttt{blk\_account\_io\_start}, \texttt{blk\_account\_io\_done})
may need adjustment on kernels that have renamed these symbols.
Second, M3 is based on $N=2$ complete fault episodes; statistically
robust MTTR characterisation requires a controlled fault-injection
campaign with many independently sampled episodes. Third, the policy
enforcement layer ran in Python fallback mode throughout, so the
isolation properties of the WASM sandbox were not validated; FSM
constraints were enforced, but the secondary enforcement boundary
did not operate with process-level isolation. Fourth, only postgres
generated qualifying stateful fault events; nginx and redis did not
produce \texttt{vfs\_write\_error} or high-latency
\texttt{blk\_io\_latency} signals during the evaluation window.
Fifth, workload classification is static: each container is assigned
to exactly one track at configuration time, and runtime
reclassification in response to changing workload characteristics
is not supported.

\subsection{Future Work}

Production deployment on Kubernetes worker nodes would replace the
per-host eBPF attachment model with a DaemonSet-based probe agent,
enabling evaluation across a realistic multi-node cluster with
heterogeneous workloads. A fault injection harness based on
\texttt{tc netem} for network faults, \texttt{stress-ng} for CPU and
I/O pressure, and filesystem namespace manipulation for mount events
would enable statistically robust M3 measurement and calibration of
the $\Lambda$, $\theta$, and $\tau$ parameters. Runtime workload
reclassification via an online classifier would eliminate the static
assignment limitation, allowing containers that acquire persistent
state dynamically (e.g., a cache that enables write-behind flushing)
to migrate from Type~S to Type~F without restart. Formal TLA+
verification of the FSM transition table in Table~\ref{tab:fsm}
would provide a machine-checked proof of safety invariants
(specifically, that \texttt{checkpoint\_and\_restart} cannot fire
from any state reachable from Healthy without traversing Audited).
Finally, installation and validation of the \texttt{wasmtime} WASM
sandbox would complete the policy enforcement layer and enable M6
to reflect true sandbox isolation rather than Python fallback enforcement.

%------------------------------------------------------------------------
\section{Conclusion}
\label{sec:conclusion}

We presented UnderCurrent, a dual-track autonomic control plane that
applies independent remediation strategies to stateless and stateful
containerised workloads. The Type~S track uses a lightweight threshold
DAG over eBPF process and network signals; the Type~F track adds a
formal five-state FSM and a policy enforcement layer that together
enforce action sequencing and irreversibility constraints. A dual
real/shadow execution model produces a divergence metric that quantifies
the safety contribution of accumulated FSM state history in every
production cycle.

Real-system evaluation on live Docker containers with active eBPF
probes confirms zero misclassification (M1 = 0.0000), perfect
explanation completeness (M5 = 1.0000), and complete reversibility of
all dispatched actions on both tracks (M4 = 1.0000). A structured
fault-injection experiment induced 44 complete fault-response cycles
in the Type~F track, yielding a divergence rate of M2 = 0.3697 and an
FSM suppression rate of M6 = 0.3697. The equality of these two
metrics reflects the regularised fault-recovery pattern and confirms
that FSM history materially changes decisions in exactly the cycles
where the safety gate is active. The nominal baseline run provides
evidence that the system does not over-trigger on healthy workloads---a
necessary precondition for operator trust in autonomous remediation.

%------------------------------------------------------------------------
\bibliographystyle{splncs04}
\begin{thebibliography}{10}

\bibitem{mapek}
Kephart, J.O., Chess, D.M.: The vision of autonomic computing.
Computer \textbf{36}(1), 41--50 (2003)

\bibitem{burnsk8s}
Burns, B., Grant, B., Oppenheimer, D., Brewer, E., Wilkes, J.:
Borg, Omega, and Kubernetes. Queue \textbf{14}(1), 70--93 (2016)

\bibitem{bcc}
Gregg, B.: BPF Performance Tools: Linux System and Application
Observability. Addison-Wesley Professional (2019)

\bibitem{cilium}
Cilium Project: eBPF-based networking, observability and security.
\url{https://cilium.io} (2024)

\bibitem{fowler}
Fowler, M.: Dark Launching.
\url{https://martinfowler.com/bliki/DarkLaunching.html} (2010)

\bibitem{opa}
Open Policy Agent: Policy-based control for cloud native environments.
\url{https://www.openpolicyagent.org} (2024)

\bibitem{tlap}
Lamport, L.: Specifying Systems: The TLA\(^+\) Language and Tools
for Hardware and Software Engineers. Addison-Wesley (2002)

\bibitem{stateright}
Creager, D.: Stateright: A model checker for implementing distributed systems.
\url{https://github.com/stateright/stateright} (2023)

\bibitem{bpftrace}
Gregg, B., Mavinakayanahalli, A.: \texttt{bpftrace} one-liners.
\url{https://github.com/bpftrace/bpftrace} (2024)

\bibitem{envoy}
Envoy Project: WebAssembly for Envoy.
\url{https://www.envoyproxy.io/docs/envoy/latest/api-v3/extensions/wasm} (2024)

\end{thebibliography}

\end{document}
